\documentclass[cs4size,a4paper,nofonts]{ctexart}
\usepackage[utf8]{inputenc}
\def\tjf{{\tt{田劲锋}}}
\def\titlec{Linux进程控制与通信}

\usepackage[a4paper,margin=2.2cm]{geometry} % 页面设置
\usepackage[unicode,breaklinks=true,
colorlinks=true,linkcolor=black,anchorcolor=black,citecolor=black,urlcolor=black,
pdftitle={\titlec},pdfauthor={\tjf}]{hyperref}
\usepackage{latexsym,amsmath,amssymb,bm}
\usepackage{graphicx}
\usepackage{xcolor}
\usepackage{fancyvrb}
\DefineShortVerb{\|}
\fvset{frame=single}

\setmainfont{Times New Roman}
\setCJKmainfont[BoldFont={SimHei}]{SimSun}  % 主要字体：宋体、黑体
\setCJKsansfont[BoldFont={STZhongsong}]{STFangsong} % 次要字体：仿宋、中宋
\setCJKmonofont{KFKai} % 等宽字体：楷体

\usepackage{listings} % 排版代码
\lstset{language=C,
  numbers=left,
  numberstyle=\small,
  basicstyle=\tt\color{black},
  commentstyle=\color{gray},
  keywordstyle=\bfseries\color{blue},
  stringstyle=\color{teal},
  showstringspaces=false,
  frame=tb,
  extendedchars=false,
%     % identifierstyle=\sf,
%     breaklines=true,
  texcl=true,
  mathescape=true,
  escapeinside=` `,
%     % texcl=true,  % 为true时，注释内容被当做LaTeX代码来处理
}


\CJKsetecglue{\hspace{0.1em}}
\renewcommand\CJKglue{\hskip -0.3pt plus 0.08\baselineskip}
\frenchspacing
\widowpenalty=10000
\linespread{1.2} % 行距

\CTEXsetup[name={实验,},number={\arabic{part}}]{part}
\CTEXsetup[name={,.},format={\large}]{section}

\usepackage[inline]{enumitem} % 调整列表样式
\setlist{noitemsep}
\setlist[itemize]{topsep=0pt,partopsep=0pt,itemsep=0pt,parsep=0pt}
\setlist[enumerate]{topsep=0pt,partopsep=0pt,itemsep=0pt,parsep=0pt}
\setlist[enumerate,1]{label={\arabic*.}}
\setlist[enumerate,2]{label={(\arabic*)}}

%\makeindex
\pagestyle{plain}

\newcommand{\image}[3][height=10cm]{%\begin{minipage}[t]{0.5\textwidth}
    \centering
    \includegraphics[#1]{#2}
    \caption{#3}
    \label{fig:#3}
%\end{minipage}
}


\begin{document}

%%%% 开始 %%%%

\setcounter{part}{1}
\begin{titlepage}

\begin{center}

\includegraphics[height=1cm]{image/haut.png}

\vspace*{1cm}
{\liti\fontsize{48pt}{50pt}{课\quad 程\quad 设\quad 计}}

\vspace*{4cm}
{\fontsize{36}{80}\sf\bfseries \titlec}

\vspace*{1cm}
{\fontsize{30}{70}\sf\bfseries \titlee}

\vfill
{\large
\newcommand{\ctline}[2]{\makebox[6em][s]{\bf #1}：\underline{\makebox[14em][c]{\qquad #2\qquad}}\\}
\ctline{课程设计名称}{数据结构课程设计}
\ctline{专业班级}{计算机 1303 班}
\ctline{学生姓名}{\tjf}
\ctline{学号}{201316920311}
\ctline{指导教师}{白\quad 浩}
\ctline{课程设计时间}{\today}
}

\end{center}

\end{titlepage}


\iffalse
\section{实验目的}
\begin{enumerate}
\item 进一步认识并发执行的概念，认识父子进程及进程创建原理；
\item 了解Linux系统中进程通信的基本原理。
\end{enumerate}

\section{实验环境}
一台装有Linux操作系统（Fedora 7），至少具有256M内存的微机。

\section{预备知识}
\begin{enumerate}
\item gcc编译器的使用
\item fork系统调用：创建一个新进程
\item getpid系统调用：获得一个进程的pid
\item wait系统调用：发出调用的进程等待子进程结束
\item pipe系统调用：建立管道
\item write系统调用：向文件中写数据
\item read系统调用：从文件中读数据
\end{enumerate}
\fi

\section{实验内容}
\begin{enumerate}[label={(\arabic*)}]
\item 编写一段程序（程序名为|parent_child.c|），使用系统调用fork()创建两个子进程，如果是父进程显示“Parent Process: A”，子进程分别显示“This is child1 (pid1 =xxxx )process: B”和“This is child1 (pid1 =xxxx )process: C”，其中“xxxx”分别指明子进程的pid号。
\item 编写一段程序(程序名为|comm.c|)，父子进程之间建立一条管道，子进程向管道中写入“Child process 1 is sending a message!”，父进程从管道中读出数据，显示在屏幕上。
\end{enumerate}

\section{实验要求}
\begin{enumerate}
\item 将parentchild.c源程序，及程序执行结果写入实验报告；
\item 将fork()系统调用后内核的工作原理写入实验报告；
\item 将comm.c源程序，及程序执行结果写入实验报告；
\item 将Linux系统中管道通信的工作原理写入实验报告。
\end{enumerate}

\section{实验步骤}

\begin{enumerate}

\item 以下是{\tt parent\_child.c}的源代码：

\lstinputlisting[caption={\tt parent\_child.c}]{exp2/parent_child.c}

首先要知道|fork()|函数定义在|<unistd.h>|中。在主函数中，9--11行尝试创建一个子进程1并将其进程ID赋值给|pid1|。12--14行判断程序在执行子进程1，打印出相应信息；注意|getpid()|会返回一个|pid_t|类型的值，应将其转换为|long|输出。15--24是父进程，其中18--20行创建子进程2，21--24行判断是子进程2打印出相应信息。
编译运行：

\begin{Verbatim}
$ cc parent_child.c -o parent_child
$ ./parent_child
Parent Process: A
This is child1 (pid1=13221) process: B
This is child2 (pid2=13222) process: C
\end{Verbatim}

注意这里的程序使用了一个小技巧，在创建了第一个子进程后，应该在父进程中再建立另外一个子进程，所以使用了嵌套的if语句。
% 如果并行的话，实际上会创建出来3个子进程，而不是2个，这不是我们所期望的。
另外一种可行的方式是，使用|wait()|来等待子进程1结束后再|fork()|出来子进程2。

\item 由|fork()|创建的是子进程。它被调用一次，返回两次。区别在于，子进程返回值是0，而父进程的返回值则是新建子进程的进程ID。子进程和父进程继续执行|fork()|调用之后的指令。子进程是父进程的副本，获得父进程的数据空间、堆和栈的副本。父进程和子进程并不共享这些存储部分，它们之间共享的是正文段。在实际实现中，操作系统通常并不完全复制其数据段和堆栈，而是采用了写时复制的技术来提高效率。

\item 以下是{\tt comm.c}的源代码：

\lstinputlisting[caption={\tt comm.c}]{exp2/comm.c}

|pipe()|也是在|<unistd.h>|中定义的。程序中，第6行定义了一个常量，表示一个文本行的最大长度。16-19行创建了一个管道，这个管道是从|fd[0]|读入，向|fd[1]|写出的。21-23行创建了一个子进程。24-26行，子进程关闭{\bf 读端}，向{\bf 写端}写出指定字符串，注意|write()|函数要求传入字符串的长度。27-31行是父进程，父进程关闭{\bf 写端}，从{\bf 读端}读取|n|个字节，写出到标准输出。这里使用了常量|STDOUT_FILENO|表示标准输出的文件号，默认是1。
编译运行：

\begin{Verbatim}
$ cc comm.c -o comm
$ ./comm
Child process 1 is sending a message!
\end{Verbatim}

这个程序同时用到了|read()|和|write()|用来读写文件描述符。

\item 管道是UNIX上最古老的进程间通信（IPC）方式。它是半双工的，即数据只能在一个方向上流动；它只能在具有公共祖先的两个进程之间使用。管道使用|pipe()|创建，返回两个文件描述符，|fd[0]|用于读，|fd[1]|用于写，|fd[1]|的输出就是|fd[0]|的输入。单个进程中的管道几乎毫无用处。通常，进程会在调用|pipo()|后接着调用|fork()|，从而创建从父进程到子进程的IPC通道。对于一个从子进程到父进程的管道，父进程关闭写端|fd[1]|，子进程关闭读端|fd[0]|。当读一个写端已关闭的管道时，在所有数据被读取后，|read()|返回0表示文件结束。当写一个读端已关闭的管道，则产生信号|SIGPIPE|，如果忽略该信号，则|write()|返回-1，并设置|errno|为|EPIPE|。


\begin{figure}[htp]
\image[width=\textwidth]{exp2/progs.png}{实际执行过程}
\end{figure}

图\ref{fig:实际执行过程}是编译和执行以上两个程序的过程。

\end{enumerate}

\end{document}
